\section{Exact translation germs and local resolution losses}
\label{sec:surgery}

A small perturbation of a multiple point can create a tiny triangle and
destroy distant old cells simultaneously. This section gives an exact
description of that wall crossing. It does not assume that smoothing
preserves or increases the triangle count.

\subsection{Simultaneous classification by affine coefficients}

Translate one indexed line parallel to itself, writing its equation as
$a_px+b_py=c_p+\varepsilon$ and leaving the others fixed. Pairwise
nonparallelism is preserved. Every determinant nondegeneracy test and
oriented half-plane test in the triangle certificate has the form
$u+v\varepsilon$: multilinearity and the single changed constant term
exclude higher powers. For positive sufficiently small $\varepsilon$,
\[
 u+v\varepsilon\ge0\quad\Longleftrightarrow\quad
 (u>0)\ \text{or}\ (u=0\text{ and }v\ge0),
\]
with analogous rules for nonpositivity and nonzero tests. A
\emph{triangle germ} is the static conjunction of these coefficient
conditions for one ordered supporting triple.

\begin{theorem}[Exact eventual translation count]\label{surgery:germs}
For any finite pairwise nonparallel arrangement and one translated line,
there is $\eta>0$ such that, simultaneously for every supporting triple
and every $0<\varepsilon<\eta$, the moved triangle predicate is equivalent
to its germ. Let $T(0)$ be the old count, and let $\mathrm{born}$ and
$\mathrm{lost}$ count germs absent from the old triangle set and old
triangles absent from the germ set, respectively. Then throughout that
interval,
\begin{equation}
 T(\varepsilon)+\mathrm{lost}=T(0)+\mathrm{born}.
 \label{surgery:birth-loss}
\end{equation}
\end{theorem}
\begin{proof}
For each affine test with nonzero constant coefficient, choose a positive
threshold preserving its sign; zero constants are decided exactly by the
linear coefficient. There are finitely many tests, supporting triples
and transverse lines. Intersect their positive neighborhoods to obtain
one simultaneous threshold. Nondegeneracy and the common half-plane
tests then give exactly the original geometric certificate predicate.
The finite set identity follows by partitioning both the old and germ
sets into their intersection and their differences. These are
\proofref{Kobon/OpenMathTranslationGerms.lean}{205}{\lean{simultaneous_triangle_stability}}
and \proofref{Kobon/OpenMathTranslationCounting.lean}{98}{\lean{translated_birth_loss_identity}}.
The lower-bound wrappers supply actual real-line witnesses from germ
counts, rather than treating a formal count as a geometric oracle.
\end{proof}

Explicit corner-incidence conditions make every old triangle's zero
side tests stationary; their other signs persist. This proves no loss
and a one-triangle gain when a new germ is present, in
\proofref{Kobon/OpenMathTranslationRetention.lean}{98}{\lean{germ_gain_one}}.
If there is exactly one concurrent supporting triple and no other line
through its center, translating one of its supports creates exactly one
tiny-triangle germ. Other new germs would have to be old degenerate
supporting triples and are excluded by uniqueness. Consequently
\begin{equation}
 T(\varepsilon)+\mathrm{lost}=T(0)+1.
 \label{surgery:isolated}
\end{equation}
This is
\proofref{Kobon/OpenMathIsolatedTripleResolution.lean}{78}{\lean{isolated_resolution_identity}},
using the local birth theorem in \lean{OpenMathTripleBirth}.
The loss set remains the \emph{global} germ difference. These theorems
use only the standard logical axioms and permit arbitrary finite order.

\subsection{A four-line counterexample to a stated half-plane yield}
\label{surgery:four}

Liu's versioned desingularization preprint, Theorem~1.1 on printed page
2, states a formula $T(A')=T(A)+1-\tau_i$, with $\tau_i$ defined by
old triangle interiors lying in the selected open half-plane
\cite{liuDesingularization2026}. Its hypotheses include one isolated
triple point and otherwise ordinary vertices. The following example
satisfies those hypotheses:
\begin{equation}
 y=0,\qquad y=x,\qquad y=-x,\qquad 2x+y=2.
 \label{surgery:four-lines}
\end{equation}
Index these supports $0,1,2,3$. The old triangular cells have supports
$013$ and $023$, and both interiors satisfy $x-y>0$.
Move support 1 to $y=x+\varepsilon$, so its motion is into $x-y<0$.
The stated half-plane count is zero. Yet for every $0<\varepsilon<1$
the moved arrangement has exactly two triangles, with supports $012$
and $023$. One old cell is lost and one is born; the proposed yield
would instead predict three triangles.

This is an interval counterexample within the formula's stated scope,
not a criticism of unrelated results in the preprint. Completeness of
the old and moved support lists, the unique triple, the whole parameter
interval, and positivity throughout the old triangle interiors are
proved in
\proofref{Kobon/OpenMathFourLineResolution.lean}{190}{\lean{halfplane_yield_fails}}.
The exact germ identity~\eqref{surgery:birth-loss} accounts for the loss
without replacing it by the half-plane count.

\subsection{Three losses and a reversed gain of two}
\label{surgery:six}

A second arrangement is
\begin{equation}
 y=0,\quad y=x,\quad y=-x,\quad x+2y=3,\quad
 -2x+y=2,\quad x+4y=-4.
 \label{surgery:six-lines}
\end{equation}
It has one concurrent triple and six old triangles. For every
$0<\varepsilon<1/100$, replacing the first line by $y=\varepsilon$
retains all six and creates one, whereas replacing it by
$y=-\varepsilon$ destroys three old triangles and creates one. The
old affected cells occupy alternating radial sectors $0,2,4$.
\begin{table}[tb]
\centering\small
\begin{tabular}{@{}lrrrr@{}}
\toprule
Motion & Old count & Lost & Born & Moved count\\
\midrule
Four-line positive resolution & 2 & 1 & 1 & 2\\
Six-line upward resolution & 6 & 0 & 1 & 7\\
Six-line downward resolution & 6 & 3 & 1 & 4\\
\bottomrule
\end{tabular}
\caption{Complete real-parameter interval counts for the displayed
arrangements. Birth and loss refer to supporting-triple sets, and the
interval conclusions are Lean theorems.}
\label{surgery:counts}
\end{table}
The complete support classification is proved in
\proofref{Kobon/OpenMathSixLineResolution.lean}{557}{\lean{up_count}}
and \proofref{Kobon/OpenMathSixLineResolution.lean}{566}{\lean{down_count}}.
The reverse collapse from the downward arrangement to the singular
one gains two triangles, in
\proofref{Kobon/OpenMathSixLineResolution.lean}{570}{\lean{reversal_gain_two}}.
Thus neither a universal cap of two destroyed cells nor a one-unit
reverse yield is valid in this scope. The positive direction also shows
that resolution can be useful when a genuine global retention proof
accompanies its local birth certificate.
